\documentclass[10pt,a4paper]{amsart}
\usepackage[margin=19mm]{geometry}
\usepackage{amsmath,amssymb,mathtools}
\usepackage[scr=rsfs,cal=euler]{mathalfa}
\usepackage{xfrac}
\usepackage[unicode,hidelinks,bookmarks=false]{hyperref}
\newcommand{\ie}{i.e.\ }
\newcommand{\cf}{cf.\ }
\newcommand{\resp}{respectively\ }
\newcommand{\Subsection}{\subsection}
\providecommand{\nobreakdash}{\nobreak\hbox{-}}
\setlength{\emergencystretch}{2em}
\multlinegap0pt
\usepackage{xcolor}
\usepackage{enumerate}
\newtheorem{prop}{Proposition}
\title{On the $p$-variation of Riemann's ``nondifferentiable'' function}
\author{Martin Lind}
\date{Source: arXiv:2608.21146v2 (CC BY 4.0)}
\begin{document}
\maketitle

\noindent\emph{Source and licence.} On the $p$-variation of Riemann's ``nondifferentiable'' function. Version arXiv:2608.21146v2 (CC BY 4.0), distributed by arXiv under the Creative Commons Attribution 4.0 International licence, http://creativecommons.org/licenses/by/4.0/, as recorded in the arXiv licence metadata for this exact version and shown on the abstract page https://arxiv.org/abs/2608.21146v2. Full licence notice: \texttt{licenses/arxiv-2608.21146-CC-BY-4.0.txt}.\par\smallskip

\noindent\textit{Editorial scope.} Counted unit: the article's prop modulusLp (source0.tex lines 275--282). The article's complete original proof of it is delivered with it (source0.tex lines 283--326, one \texttt{proof} environment of 44 lines). No mathematics is written, repaired or re-derived: the statement and the proof are the article's own lines, in the article's own notation, and no retained line is substituted or re-flowed.  No further source-local context is needed: every cross-reference of every retained line resolves inside the two delivered spans.  Nothing was omitted from any retained span, so the package carries no omission record.  No retained line contains a citation, so the wrapper carries no bibliography and no bibliography entry.  No retained line contains a figure environment, an image-inclusion command or drawing code, so no image is delivered and the package declares no asset dependency.  Editorial formatting only, in addition: an AMS \texttt{amsart} preamble that restates the article's own declarations and macros used by the retained lines, namely the environments \texttt{prop} on separate counters, together with the standard packages those lines need.  The retained environments print in this document as Proposition 1 in order of appearance, whereas the article prints the same items with the numbering of its own full document.  The complete native source of the article as distributed in the arXiv e-print for this exact version is bundled unchanged as \texttt{sources/208228/source0.tex}.
\begin{prop}
    \label{modulusLp}
    For $1\le p\le2$ and any $\delta\in(0,1]$, there holds
    \begin{equation}
        \label{modulusIneq}
        \omega(\Phi;\delta)_p\le 6\delta^{3/4}.
    \end{equation}
\end{prop}

\begin{proof}
We treat first the case $p=2$. Let $h\in(0,1]$ be arbitrary. For any $x\in\mathbb{R}$, there holds
\begin{equation}
    \nonumber
    \Phi(x+h)-\Phi(x)=\sum_{n=1}^\infty(e^{\pi i n^2h}-1)\frac{e^{\pi i n^2 x}}{\pi in^2}.
\end{equation}
Applying Parseval's identity
\begin{equation}
    \label{l2norm}
    \|\Delta_h\Phi\|^2_{L^2(0,2)}=2\sum_{n=1}^\infty\frac{|e^{\pi i n^2h}-1|^2}{\pi^2 n^4}.
\end{equation}
Note that $|e^{iz}-1|\le\min(2,|z|)$, whence
\begin{equation}
    \label{elemIneq}
    |e^{\pi i n^2h}-1|\le\min(2,\pi n^2h).
\end{equation}
Let $N=\lfloor h^{-1/2}\rfloor$, use (\ref{elemIneq}), and split the sum at the right-hand side of (\ref{l2norm}) as follows:
\begin{eqnarray}
    \nonumber
    \sum_{n=1}^\infty\frac{|e^{\pi i n^2h}-1|^2}{\pi^2 n^4}&\le&\sum_{n=1}^\infty\frac{\min(2,\pi n^2h)^2}{\pi^2n^4}\\
    \nonumber
    &\le &\sum_{n\le N}\frac{\pi^2 n^4h^2}{\pi^2 n^4}+\sum_{n> N}\frac{4}{\pi^2 n^4}\\
    \nonumber
    &=&h^2\sum_{n\le N}1+\frac{4}{\pi^2}\sum_{n> N}\frac{1}{n^4}\\
    \nonumber
    &=& h^2N+\frac{4}{3\pi^2N^3}.
\end{eqnarray}
Since $N\le h^{-1/2}$, we get $h^2N\le h^{3/2}$. Furthermore, since $h^{-1/2}\ge1$ and $N=\lfloor h^{-1/2}\rfloor$, we have $N\ge h^{-1/2}/2$. It follows that $1/N^3\le 8h^{3/2}$ and consequently $4/(3\pi^2 N^3)\le (32/3\pi^2)h^{3/2}$. Combining the preceding estimates yields
\begin{equation}
    \nonumber
    \|\Delta_h\Phi\|^2_{L^2(0,2)}\le 2\left(1+\frac{32}{3\pi^2}\right)h^{3/2}.
\end{equation}
Hence, $\omega(\Phi;\delta)_2\le3\delta^{3/4}$ for any $\delta\in(0,1]$. Furthermore, for $1\le p\le 2$
\begin{equation}
    \nonumber
    \|f\|_{L^p(0,2)}\le 2^{1/p-1/2}\|f\|_{L^2(0,2)}.
\end{equation}
Hence, for $1\le p\le 2$ and $\delta\in(0,1]$
\begin{equation}
    \nonumber
    \omega(\Phi;\delta)_p\le5\delta^{3/4},
\end{equation}
thus proving (\ref{modulusIneq}).
\end{proof}

\end{document}
